\abstracten

Structuring course knowledge supports teaching and research in computer networking, yet such documents are dominated by specialized terms and cross-chapter relations, making manual knowledge graph construction costly and hard to scale. This work automates knowledge graph construction from unstructured course materials with large language models (LLMs).

A six-stage pipeline converts PDFs into Markdown, text chunks, raw triples, a merged knowledge graph, and Neo4j storage. Multimodal vision-language models perform OCR (Optical Character Recognition) and produce structured Markdown; a recursive header-aware chunking algorithm preserves document hierarchy within token limits. Chinese prompts extract four entity types (protocol, concept, mechanism, metric) and four relation types (contains, depends\_on, compared\_with, applied\_to) under zero-shot and few-shot settings, with an evidence verification gate grounding each triple in source text. Entity canonicalization combines FAISS (Facebook AI Similarity Search) with Union-Find to merge duplicates.

The pipeline processed 205 course documents and generated 16,438 chunks ($\sim$10.4 million tokens). The final knowledge graph (\texttt{merged-fs-recheck}, with LLM-based candidate review enabled) contains 32,760 entities and 17,627 relational triples. On 500 expert-reviewed gold triples, a four-tier matching protocol (Strict / Substring / Slot / Schema-level) yields triple-level F1 of 15\%--19\% under substring, 37\%--42\% under slot, and 74\%--82\% under schema-level matching. Frequent entities such as TCP, UDP, HTTP, and IP indicate that the graph captures core networking concepts, providing a structured knowledge base for networking education.

\vspace{\baselineskip}

\keywordsen Large Language Models, Named Entity Recognition, Relation Extraction, Knowledge Graph, Course Knowledge, Prompt Engineering, Entity Canonicalization
